% TOP_PROBLEM: {"id": 10000038, "title": "Cheeger constant and the percolation nonuniqueness phase", "category_id": 19, "category": {"id": 19, "name": "probability", "display_name": "Probability", "description": "Problems involving probability theory, stochastic processes, and random structures.", "slug": "probability", "order_index": 19, "created_at": "2026-07-31T15:26:25.671Z"}, "status": "partially_solved", "problem_number": "AMR-099-0038", "set_id": 13}
% FIRST_POSTED: 2026-10-06
% ATTEMPT_STATUS: solved
% OPENAI_CLAIM: {"schema_version":1,"source_repo":"https://github.com/openai/math","source_commit":"adc7f1241b42e322a6451854ab7e4b4c146bf78a","release_date":"2026-10-06","match":"stronger","resolution":"proved","independently_reviewed":false,"second_pass":"agreed","families":[{"family":"214","title":"The Benjamini–Schramm nonuniqueness conjecture","subject":"Probability and statistical mechanics","manuscripts":[{"title":"Nonuniqueness of percolation on nonamenable quasi-transitive graphs","pdf_path":"preprints/Nonuniqueness-of-percolation-on-nonamenable-quasi-transitive-graphs-September-24-2026/paper.pdf","readme_path":"preprints/Nonuniqueness-of-percolation-on-nonamenable-quasi-transitive-graphs-September-24-2026/README.md","date":"2026-09-24","theorem_ref":"thm:main (sections/introduction.tex:58)","statement_sources":[{"path":"preprints/Nonuniqueness-of-percolation-on-nonamenable-quasi-transitive-graphs-September-24-2026/build/sections/introduction.tex","line":58,"label":"thm:main"}]}],"lean":{"doc_path":"lean/docs/214.md","comparators":["lean/ComparatorChallenges/BenjaminiSchramm.json","lean/ComparatorChallenges/BenjaminiSchramm.lean","lean/ComparatorChallenges/CayleyPercolation.json","lean/ComparatorChallenges/CayleyPercolation.lean"],"declarations":[],"covers_main_theorem":true},"reasoning_trace":null}]}
\documentclass[11pt]{article}
\usepackage[margin=25mm]{geometry}
\usepackage{fontspec}
\setmainfont{Latin Modern Roman}
\usepackage{amsmath,amssymb,mathtools}
\usepackage{unicode-math}
\setmathfont{Latin Modern Math}
\usepackage{enumitem}
\usepackage{xurl,hyperref,bookmark}
\hypersetup{colorlinks=true,urlcolor=blue}
\setcounter{secnumdepth}{0}
\setlength{\parindent}{0pt}
\setlength{\parskip}{5pt}
\setlength{\emergencystretch}{3em}
\begin{document}
\section*{Cheeger constant and the percolation nonuniqueness phase}

\subsection{OpenAI's claim}
OpenAI presents a proof in the selected manuscript. Every nonamenable locally finite quasi-transitive graph is claimed to satisfy $p_c<p_{2\to2}\leq p_u$, yielding a nonempty phase with infinitely many infinite clusters.

\subsection{Mathematical statement of the claimed result}
\paragraph{Result 1.} Let \(G=(V,E)\) have infinitely many vertices, finite degrees, connectedness and quasi-transitivity, with positive vertex Cheeger constant
\[
 h_V(G)=\inf_{0<|A|<\infty}|\partial_VA|/|A|>0.
\]
For independent bond percolation, let \(T_p(x,y)=\mathbb P_p(x\leftrightarrow y)\), acting on \(\ell^2(V)\) with counting measure, and \(p_2=\sup\{p:\|T_p\|_{2\to2}<\infty\}\). With \(p_c\) and \(p_u\) the appearance and uniqueness thresholds, the source claims
\[
 \sup_{p<p_c}\|T_p\|_{2\to2}=\|T_{p_c}\|_{2\to2}<\infty,
 \qquad p_c<p_2\le p_u.
\]
It further supplies \(p_c<p_1<p_2'<1\) such that the uniform-edge-label coupling almost surely has infinitely many infinite clusters simultaneously at every \(p\in[p_1,p_2']\).

\subsection{Outline of the argument}
For the new unimodular case, local endpoint maps control moments of finite critical clusters. A deformation lowers susceptibility while retaining a large connectivity-operator norm, thereby proving the strict operator-threshold separation; the nonunimodular case uses the established operator result.

\subsection{Scope relative to this problem}
The cited theorem proves pc<p2to2<=pu for every nonamenable locally finite quasi-transitive graph, hence the nonamenable direction for all vertex-transitive graphs in our equivalence; the reverse is the established uniqueness theorem for amenable transitive Bernoulli percolation. The claimed conclusion supplies the requested case and extends it in the directions just specified.

\subsection{Formal verification}
For family 214, the documented statements include the strict operator-threshold separation, the nonunique percolation interval, and selected critical moment and tail bounds, with a Cayley-graph specialization. The listed Comparator files are \texttt{BenjaminiSchramm.lean}, \texttt{CayleyPercolation.lean}. The inventory does not supply a named declaration for these Comparator entries. This site has not executed the Lean code. The scope descriptions do not certify the remaining manuscript arguments.

\subsection{Status caveat}
OpenAI attributes these manuscripts to work by an internal model. Its release cautions that arguments without formalization can contain mistakes. This catalogue has not independently checked the proofs; the labels report the mathematical claims attributed to OpenAI.

\subsection{Sources}
\begin{itemize}
\item OpenAI, \emph{Nonuniqueness of percolation on nonamenable quasi-transitive graphs}, 2026-09-24; citation key \texttt{OAI:Nonuniqueness-of-percolation-on-nonamenable-quasi-transitive-graphs-September-24-2026}. \href{https://github.com/openai/math/blob/adc7f1241b42e322a6451854ab7e4b4c146bf78a/preprints/Nonuniqueness-of-percolation-on-nonamenable-quasi-transitive-graphs-September-24-2026/paper.pdf}{PDF (GitHub)}; \href{https://github.com/openai/math/raw/adc7f1241b42e322a6451854ab7e4b4c146bf78a/preprints/Nonuniqueness-of-percolation-on-nonamenable-quasi-transitive-graphs-September-24-2026/paper.pdf}{PDF (direct)}; \href{https://github.com/openai/math/blob/adc7f1241b42e322a6451854ab7e4b4c146bf78a/preprints/Nonuniqueness-of-percolation-on-nonamenable-quasi-transitive-graphs-September-24-2026/README.md}{Manuscript page and citation}; \href{https://github.com/openai/math/blob/main/preprints/Nonuniqueness-of-percolation-on-nonamenable-quasi-transitive-graphs-September-24-2026/paper.pdf}{Latest version}; Statement source: \href{https://github.com/openai/math/blob/adc7f1241b42e322a6451854ab7e4b4c146bf78a/preprints/Nonuniqueness-of-percolation-on-nonamenable-quasi-transitive-graphs-September-24-2026/build/sections/introduction.tex\#L58}{thm:main}.
\item Family 214: \href{https://github.com/openai/math/blob/adc7f1241b42e322a6451854ab7e4b4c146bf78a/lean/docs/214.md}{Lean scope}.
\item \href{https://github.com/openai/math/blob/adc7f1241b42e322a6451854ab7e4b4c146bf78a/lean/ComparatorChallenges/BenjaminiSchramm.json}{Comparator BenjaminiSchramm.json}.
\item \href{https://github.com/openai/math/blob/adc7f1241b42e322a6451854ab7e4b4c146bf78a/lean/ComparatorChallenges/BenjaminiSchramm.lean}{Comparator BenjaminiSchramm.lean}.
\item \href{https://github.com/openai/math/blob/adc7f1241b42e322a6451854ab7e4b4c146bf78a/lean/ComparatorChallenges/CayleyPercolation.json}{Comparator CayleyPercolation.json}.
\item \href{https://github.com/openai/math/blob/adc7f1241b42e322a6451854ab7e4b4c146bf78a/lean/ComparatorChallenges/CayleyPercolation.lean}{Comparator CayleyPercolation.lean}.
\item \href{https://github.com/openai/math/blob/adc7f1241b42e322a6451854ab7e4b4c146bf78a/CONTENTS.md}{OpenAI catalogue entry}.
\end{itemize}
\end{document}
